\section{결과 및 논의}
\label{sec:discussion}

\subsection{기존 도구·연구와의 기능 비교}

\begin{table*}[t]
\centering\small
\caption{정간보 관련 도구·연구와의 기능 비교
(\(\bullet\)=지원, \(\circ\)=부분 지원, ---=미지원)}
\label{tab:compare}
\begin{tabular}{@{}lcccccc@{}}
\toprule
 & 본 시스템 & XML 기술~\cite{lee2013xml} & 웹 설계~\cite{jkca2022web}
 & 웹 편집기~\cite{jeongganbo-editor-web} & HWP 표 조판 & MuseScore~\cite{musescore} \\
\midrule
사람이 읽는 텍스트 문법      & \(\bullet\) & --- & --- & --- & --- & --- \\
WYSIWYG 편집                & \(\bullet\) & --- & \(\circ\) & \(\bullet\) & \(\bullet\) & \(\bullet\) \\
분박·붙임(위치 시가)         & \(\bullet\) & \(\bullet\) & \TODO{확인} & \(\circ\) & \(\circ\) & --- \\
시김새                      & \(\bullet\)(59종) & \(\circ\) & \TODO{확인} & --- & \(\circ\)(수작업) & --- \\
장단·가사 계층               & \(\bullet\) & \(\circ\) & \TODO{확인} & --- & \(\circ\)(수작업) & --- \\
우종서·대강 조판             & \(\bullet\) & \(\circ\) & \TODO{확인} & \(\circ\) & \(\circ\)(수작업) & --- \\
재생                        & \(\bullet\) & \(\circ\) & --- & --- & --- & \(\bullet\)(오선보) \\
인쇄·PDF·PNG 출력            & \(\bullet\) & --- & --- & \(\circ\)(PNG) & \(\bullet\) & \(\bullet\) \\
구조화 문서 교환             & \(\bullet\)(JSON) & \(\bullet\)(XML) & \TODO{확인} & \(\circ\)(JSON) & --- & \(\bullet\)(MusicXML) \\
서버리스·오프라인            & \(\bullet\) & 해당 없음 & \TODO{확인} & --- & \(\bullet\) & \(\bullet\) \\
\bottomrule
\end{tabular}
\end{table*}

표~\ref{tab:compare}에 정간보 관련 도구·연구와의 기능 비교를 정리했다.
본 시스템은 텍스트 문법과 WYSIWYG을 겸비하고 시김새·장단·가사·조판·출력·
재생을 한 도구에서 지원한다는 점에서, 인코딩 제안에 머문 연구나 기본
격자 입력만 제공하는 기존 편집기와 구별된다.
\TODO{비교표의 각 칸을 해당 도구를 직접 실행/원문 확인 후 재검증할 것 ---
특히 jkca2022web 열은 현재 추정임}

\subsection{사례 연구: 간행 악보 재현}

본문 그림들(그림~\ref{fig:ui}--\ref{fig:mapping})의 문서는 문법 요소를
고루 포함하도록 구성한 예시 악보이다. 간행 악보 재현의 사례 연구로는
해금정악보의 군악(12정간, 대강 3·3·3·3)을 대상으로 원본 지면과 동일한
배치(세로 제목 칸·장단 열·시김새 포함)를 재현하고 원본 스캔과 나란히
비교할 계획이다.
\TODO{저자 직접 작업 필요: `악보 예시/예시 이미지-12박(대강3333).png'
(해금정악보 68쪽 군악)를 채보해 재현하고 --- 해금 활대 방향 표시 등
관악보 고유 기호의 처리 방식 결정 포함 --- 원본/출력 비교 그림, 입력
소요 시간, 텍스트 소스 길이, \code{.jgb.json} 파일 크기 등 정량 정보를
기입할 것. 원본 스캔 게재 시 저작권 확인 필요}
전 과정에 서버 통신이 없고, 완성 문서는 수십 KB 수준의
\code{.jgb.json} 하나로 보존·공유된다.

\subsection{한계}

현재 구현의 한계는 다음과 같다.
첫째, 재생이 사인파 단선율에 그치며 시김새·장단이 소리에 반영되지 않는다.
시김새는 본질적으로 미분음적 음높이 굴곡이므로, 기호별 음고 궤적
모델링이 필요하다.
둘째, 셀 서식(배경색·테두리)이 선율 정간에만 적용되고 장단·가사 칸은
지원하지 않는다.
셋째, MusicXML·MEI 등 표준 포맷과의 상호 변환이 없어 오선보 병기나
서양 도구 연계가 수동이다.
넷째, 사용성 평가를 아직 수행하지 않았다. 국악 전공자·교사를 대상으로
과제 기반 평가(기존 HWP 조판 대비 작성 시간·오류율·SUS 점수)를 수행하는
것이 후속 과제다. \TODO{투고 전 소규모라도 사용자 평가를 실시해 5장을
평가 결과 중심으로 재구성하는 것을 강력 권장 --- JDCS 심사에서 평가 부재가
지적될 가능성이 높음}
